\section{托卡马克：用磁场装住太阳}

上一章讲反应原理，本章进入主流技术路线。地球不能像太阳一样靠巨大重力压住聚变反应，因此人类主要尝试两类约束：惯性约束和磁约束。惯性约束用极短强激光等方式压缩燃料靶丸，适合研究高能量密度物理和武器相关问题，但不适合稳定民用发电；磁约束则用磁场把高温等离子体约束在空间中，避免它直接接触固体壁。

托卡马克是磁约束里最主流、实验参数最高的路线。它的直觉是一个甜甜圈形磁场：让等离子体沿环形路径运动，在环向磁场和环向电流共同形成的位形中被约束。访谈中杨钊强调，托卡马克之所以成为共识，不是因为名字好听，而是因为过去几十年里它在实验参数上明显领先其他磁约束路线。

\begin{figure}[H]
\centering
\includegraphics[width=0.92\textwidth]{tokamak-stack.png}
\caption{托卡马克系统栈：磁体、真空室、等离子体控制和能量转换层层耦合。自制概念图，依据 00:06:32--00:16:00 与 01:00:50--01:04:36 对谈内容整理。}
\end{figure}

\begin{knowledgebox}{读图：托卡马克是系统工程，不是单个磁铁}
图中从超导磁体到真空室、加热系统、控制系统和热工系统，说明托卡马克是一台多系统耦合机器。磁体提供约束，真空室提供运行环境，加热和控制把等离子体推到目标参数，热工系统则决定能否成为电站。
\end{knowledgebox}

\subsection{铜、低温超导与高温超导}

本节看磁体材料如何改变整台装置。早期托卡马克常用铜磁体，工程简单，但铜有电阻，通大电流会发热，只能短时间脉冲运行。低温超导可以长时间运行，但需要极低温和复杂真空/冷屏系统；更关键的是低温超导临界磁场有限，要达到足够能量增益，装置就必须做得很大，典型代表是 ITER 这种超大国际项目。

高温超导的核心价值不是“温度真的很高”，而是在低温运行条件下可以承受更高磁场。磁场越强，同等性能所需装置尺寸越小；尺寸越小，建造成本和周期越低。杨钊给出的路线判断是：高温超导可能把低温超导路线中百亿欧元级、数十年周期的大装置，压缩成更适合创业公司迭代的工程对象。

\begin{knowledgebox}{术语消化：三代托卡马克磁体}
\begin{tabularx}{\textwidth}{p{0.18\textwidth}p{0.34\textwidth}X}
\toprule
阶段 & 优点 & 主要问题 \\
\midrule
铜磁体 & 工程复杂度低，适合早期实验 & 发热大，只能短脉冲运行。 \\
低温超导 & 可长时间运行，有 EAST、KSTAR、JT60SA 等装置经验 & 低温系统复杂，临界磁场限制导致装置巨大。 \\
高温超导 & 临界磁场更高，有望显著缩小装置和成本 & 材料加工、结构受力、低温集成和系统验证仍很新。 \\
\bottomrule
\end{tabularx}
\end{knowledgebox}

\begin{importantbox}{高温超导的商业化逻辑}
高温超导不是让聚变反应本身变简单，而是改变装置的工程经济性：更高磁场允许更小装置，更小装置带来更低成本、更短迭代周期和更适合创业公司推进的研发路径。
\end{importantbox}

\subsection{本章小结}

托卡马克是目前最有实验积累的磁约束路线。高温超导托卡马克的关键，不只是科学性能，而是能否用更高磁场把装置规模和度电成本压下来。

